\section{Verification and results}
\label{sec:risultati}

\subsection{Verification procedure}
\label{sec:metodo}

All test cases from the reference study were reproduced using the corresponding
zero-dimensional verification programs. The solver was also applied to the
pure-nitrogen cases and to the molecular-atomic nitrogen mixture, both without
reactions and with reactions initiated from a thermally non-equilibrium state.
These solver calculations used a single-cell mesh, with the same initial
conditions, gas data and time step of one nanosecond as the zero-dimensional
programs. The verification comprises three comparisons:

\begin{enumerate}[leftmargin=1.8em]
  \item \textbf{solver against zero-dimensional programs}: both
    implementations use the same physical model and source-term functions;
    this comparison assesses the coupling and numerical integration of the
    model within the solver;
  \item \textbf{zero-dimensional programs against the reference study}:
    this comparison assesses consistency with the published thermochemical
    results;
  \item \textbf{final temperature against energy conservation}: the
    equilibrium temperature in each non-reacting case is compared with the
    value calculated from the initial energy and with the value reported in
    the reference study.
\end{enumerate}

In addition, the modelling choices that the reference study does not specify
are examined by repeating the cases they affect with the alternative setting
(Section~\ref{sec:risultati-sensibilita}).

The \hyfoam{} curves were extracted directly from the vector figures in the
PDF of the reference study. Each curve is stored as a sequence of points,
which were converted to physical coordinates using the axis ticks. The
estimated extraction uncertainty is comparable to the line width and
corresponds to approximately 1\% in temperature. Unless otherwise stated,
temperature and species-density deviations are reported as percentages and
denote the maximum over the sampled transient. The programs record results
every 10 or 100~ns, and the comparisons use these output samples. In the
figures, light solid lines denote the reference results, dashed lines the
zero-dimensional results and dots the solver results. The figures retain the
notation $T_v$ for vibrational temperature; when electronic excitation is
included, this denotes the common vibro-electronic temperature $\Tve$.

\subsection{Pure nitrogen}
\label{sec:risultati-azoto}

The first two cases consider nitrogen at atmospheric pressure. The initial
vibrational temperature is either much lower than the translational
temperature (vibrational heating) or much higher (vibrational cooling). In
the heating case, the final temperature agrees with the value stated in the
reference study, while the maximum transient deviations are 0.3\% for the
translational temperature and 1.5\% for the vibrational temperature. In the
cooling case, the final temperature differs by 0.03\%; the corresponding
transient deviations are 2.8\% and 5.1\%, respectively, indicating slightly
delayed relaxation
(Figure~\ref{fig:azoto-3}).

\begin{figure}[ht]
  \centering
  \begin{subfigure}[t]{0.49\textwidth}
    \includegraphics[width=\linewidth]{nitrogen-heating}
    \caption{Vibrational heating.}
  \end{subfigure}\hfill
  \begin{subfigure}[t]{0.49\textwidth}
    \includegraphics[width=\linewidth]{nitrogen-cooling}
    \caption{Vibrational cooling.}
  \end{subfigure}
  \caption{Pure nitrogen at atmospheric pressure: time evolution of the
  translational ($T_{tr}$) and vibrational ($T_v$) temperatures.}
  \label{fig:azoto-3}
\end{figure}

The third case uses a substantially higher initial translational temperature,
at which electronic excitation makes a non-negligible contribution to the
energy. Calculations are performed both with and without this contribution
(Figure~\ref{fig:azoto-4}). In both configurations, the final temperature
agrees with the reference value to within 0.3\%. Including electronic energy
reduces the final temperature by approximately 20\%. The maximum transient
deviations in the translational and vibro-electronic temperatures are, respectively,
0.05\% and 0.8\% without electronic energy, and 0.3\% and 4\% with it.

\begin{figure}[ht]
  \centering
  \begin{subfigure}[t]{0.49\textwidth}
    \includegraphics[width=\linewidth]{nitrogen-hot-no-electronic}
    \caption{Without electronic energy.}
  \end{subfigure}\hfill
  \begin{subfigure}[t]{0.49\textwidth}
    \includegraphics[width=\linewidth]{nitrogen-hot-electronic}
    \caption{With electronic energy.}
  \end{subfigure}
  \caption{Pure nitrogen at high temperature. The horizontal dotted line
  indicates the equilibrium temperature calculated from energy conservation.}
  \label{fig:azoto-4}
\end{figure}

\subsection{Non-reacting mixtures}
\label{sec:risultati-miscele}

The first case considers an equimolar mixture of molecular and atomic
nitrogen, initially at a very high translational temperature and a low
vibrational temperature (Figure~\ref{fig:azoto-atomico}). The final
temperature agrees with the reference value to within 0.2\%, and the maximum
translational-temperature deviation is 2.5\%. The rise in vibrational
temperature is delayed relative to the reference result, with a deviation of
up to 20\% in the steepest portion of the curve. The most likely explanation
is the different method used by \mpp{} to average relaxation times in
mixtures. This difference is consistent with slower vibrational relaxation
while preserving agreement in the final state.

\begin{figure}[ht]
  \centering
  \includegraphics[width=0.62\textwidth]{nitrogen-atomic}
  \caption{Equimolar mixture of molecular and atomic nitrogen. The horizontal
  dotted line indicates the equilibrium temperature calculated from energy
  conservation.}
  \label{fig:azoto-atomico}
\end{figure}

The second case considers an equimolar nitrogen-oxygen mixture at atmospheric
pressure, with an initial vibrational temperature much higher than the
translational temperature. A separate vibrational temperature is assigned to
each molecular species, and the case is computed with and without V--V
exchange between them (Figure~\ref{fig:azoto-ossigeno}). In both
configurations the final temperature agrees with the reference value to
within 0.3\%. Without V--V exchange, the maximum deviations are 2.7\% for the
translational temperature, 2.7\% for the nitrogen vibrational temperature and
6.4\% for the oxygen vibrational temperature; with V--V exchange they
decrease to 2.0\%, 2.5\% and 3.4\%, respectively. Consistent with the
reference study, V--V exchange brings the two vibrational temperatures closer
to each other and accelerates relaxation towards equilibrium.

\begin{figure}[ht]
  \centering
  \begin{subfigure}[t]{0.49\textwidth}
    \includegraphics[width=\linewidth]{nitrogen-oxygen-no-vv}
    \caption{Without V--V exchange.}
  \end{subfigure}\hfill
  \begin{subfigure}[t]{0.49\textwidth}
    \includegraphics[width=\linewidth]{nitrogen-oxygen-vv}
    \caption{With V--V exchange.}
  \end{subfigure}
  \caption{Equimolar nitrogen-oxygen mixture at atmospheric pressure. The
  horizontal dotted line indicates the equilibrium temperature calculated from
  energy conservation.}
  \label{fig:azoto-ossigeno}
\end{figure}

\subsection{Reacting mixtures}
\label{sec:risultati-reagenti}

The molecular-atomic nitrogen mixture was also simulated with nitrogen
dissociation, using two initial states. The first is the thermally
non-equilibrium state of the non-reacting case
(Figure~\ref{fig:reagenti-noneq}). At the end of the simulation both
temperatures agree with the reference values to within 0.4\%, while the
species densities agree to within 1.5\% throughout the transient. The maximum
translational-temperature deviation is 2.3\%. The vibrational temperature
exhibits the same delay as in the non-reacting mixture, with deviations of up
to 17\% at early times.

\begin{figure}[ht]
  \centering
  \begin{subfigure}[t]{0.49\textwidth}
    \includegraphics[width=\linewidth]{reacting-noneq-temperatures}
    \caption{Temperatures.}
  \end{subfigure}\hfill
  \begin{subfigure}[t]{0.49\textwidth}
    \includegraphics[width=\linewidth]{reacting-noneq-densities}
    \caption{Species number densities normalised by the initial value.}
  \end{subfigure}
  \caption{Reacting molecular-atomic nitrogen mixture initialised in thermal
  non-equilibrium, with Park exponent $a = 0.7$.}
  \label{fig:reagenti-noneq}
\end{figure}

The second case starts from equal temperatures, corresponding to thermal
equilibrium but chemical non-equilibrium (Figure~\ref{fig:reagenti-eq}).
Dissociation begins immediately and cools the gas, and the vibrational
temperature decreases more slowly than the translational one. The final
temperatures agree with the reference values to within 0.8\%, and the species
densities to within 2.3\%. The maximum transient deviations in the
translational and vibrational temperatures are 2.0\% and 10\%, respectively.
The vibrational temperature again remains above the reference curve, as in
the other mixtures containing atomic nitrogen.

\begin{figure}[ht]
  \centering
  \begin{subfigure}[t]{0.49\textwidth}
    \includegraphics[width=\linewidth]{reacting-eq-temperatures}
    \caption{Temperatures.}
  \end{subfigure}\hfill
  \begin{subfigure}[t]{0.49\textwidth}
    \includegraphics[width=\linewidth]{reacting-eq-densities}
    \caption{Species number densities normalised by the initial value.}
  \end{subfigure}
  \caption{Reacting molecular-atomic nitrogen mixture initialised in thermal
  equilibrium at the same high temperature, with Park exponent $a = 0.7$.}
  \label{fig:reagenti-eq}
\end{figure}

\subsection{Five-species air}
\label{sec:risultati-aria}

The final case considers air, initially composed of 79\% nitrogen and 21\%
oxygen by mole, at low pressure and high temperature. The calculation uses
nineteen reactions and a single temperature. With the Quantum-Kinetic rate
constants, the temperature agrees with the reference result to within 0.7\%
during the transient and 0.4\% at the final time. The densities of all five
species agree to within 15\%, with four species agreeing to within 7\%
(Figure~\ref{fig:aria}).

\begin{figure}[ht]
  \centering
  \begin{subfigure}[t]{0.49\textwidth}
    \includegraphics[width=\linewidth]{air-temperature}
    \caption{Temperature.}
  \end{subfigure}\hfill
  \begin{subfigure}[t]{0.49\textwidth}
    \includegraphics[width=\linewidth]{air-densities}
    \caption{Species number densities normalised by the initial value.}
    \label{fig:aria-densita}
  \end{subfigure}
  \caption{Five-species air with the Quantum-Kinetic rate constants.}
  \label{fig:aria}
\end{figure}

\subsection{Sensitivity to the modelling choices}
\label{sec:risultati-sensibilita}

Four of the settings used above are not stated in the reference study and were
selected by comparison with its curves. Table~\ref{tab:sensibilita} collects
them, together with the alternative that was discarded and the deviation it
produces. Each adopted setting reproduces the reference results within the
deviations reported in the preceding subsections, whereas each alternative
produces deviations at least twice as large, and in most cases much larger.

The two choices that concern the chemistry are illustrated in
Figure~\ref{fig:sensibilita}. The exponent in Park's temperature determines
how strongly a low vibrational excitation slows dissociation: with the value
fixed in \mpp{} the mixture begins to dissociate earlier than in the reference
study, although the final composition is unaffected. The rate constants for
air act throughout the transient and also change the final densities. The
four choices therefore separate into two groups, which helps to identify
their effects: the treatment of electronic energy and the rate constants for
air alter the final state, and are thus constrained by the equilibrium
temperature and the final composition, whereas the energy driving the V--T
exchange and the exponent in Park's temperature alter only the path towards
that state.

\begin{table}[H]
  \centering
  \caption{Modelling choices not specified in the reference study, and the
  effect of the discarded alternative on the maximum deviation from the
  reference results.}
  \label{tab:sensibilita}
  \small
  \begin{tabular}{@{}>{\raggedright\arraybackslash}p{3.0cm}>{\raggedright\arraybackslash}p{2.2cm}>{\raggedright\arraybackslash}p{1.9cm}>{\raggedright\arraybackslash}p{7.1cm}@{}}
    \toprule
    Choice & Adopted & Alternative & Effect of the alternative \\
    \midrule
    Electronic energy in the mixtures & ground level only & all levels &
      final temperature 25\% below the reference value for molecular and
      atomic nitrogen and 59\% above it for nitrogen and oxygen, instead of
      0.3\% \\[2pt]
    \midrule
    Energy driving the V--T exchange & vibro-electronic & vibrational only &
      in high-temperature nitrogen, transient deviations of 28\% in $\Ttr$
      and 56\% in $\Tve$, instead of 0.3\% and 4\% \\[2pt]
    \midrule
    Exponent in Park's temperature & $a = 0.7$ & $a = 0.5$ &
      in the reacting mixture, species densities deviate by up to 9\% instead
      of 1.5\% and $\Ttr$ by 6.7\% instead of 2.3\% \\[2pt]
    \midrule
    Rate constants for air & Quantum-Kinetic & Park &
      temperature deviates by up to 8.9\% instead of 0.7\%, and species
      densities by factors up to 3 instead of 15\% \\
    \bottomrule
  \end{tabular}
\end{table}

\begin{figure}[ht]
  \centering
  \begin{subfigure}[t]{0.49\textwidth}
    \includegraphics[width=\linewidth]{park-exponent-densities}
    \caption{Reacting nitrogen mixture with the two exponents in Park's
    temperature.}
  \end{subfigure}\hfill
  \begin{subfigure}[t]{0.49\textwidth}
    \includegraphics[width=\linewidth]{air-densities-park-rates}
    \caption{Air with the Park rate constants, to be compared with
    Figure~\ref{fig:aria-densita}.}
  \end{subfigure}
  \caption{Effect on the species densities of the two modelling choices that
  concern the chemistry; in (a) the dotted lines are the results obtained with
  the alternative exponent.}
  \label{fig:sensibilita}
\end{figure}

\subsection{Solver versus zero-dimensional programs}
\label{sec:risultati-solver}

Table~\ref{tab:solver} summarises the maximum deviations between the solver and
the zero-dimensional programs for all cases evaluated with both
implementations. The final temperatures coincide, and the transient
temperature profiles agree to within 0.3\% in every case except
high-temperature nitrogen. In that case, the deviation reaches 1.2\% during
the rapid relaxation in the first 100~ns. This regime is more sensitive to
the time-integration procedure: the solver performs two iterations per time
step, whereas the zero-dimensional programs perform one. In the reacting
case, the species densities agree to within 0.03\%. These results support the
consistency of the solver coupling and numerical integration, including the
chemical source terms, for the cases examined.

\begin{table}[H]
  \centering
  \caption{Maximum deviation between the solver and the zero-dimensional
  programs.}
  \label{tab:solver}
  \small
  \begin{tabular}{lccc}
    \toprule
    Case & $T_{tr}$ & $T_{ve}$ & Composition \\
    \midrule
    Nitrogen, vibrational cooling                    & 0.00\% & 0.00\% & -- \\
    Nitrogen, vibrational heating                    & 0.03\% & 0.21\% & -- \\
    High-temperature nitrogen, without electronic energy & 0.14\% & 0.69\% & -- \\
    High-temperature nitrogen, with electronic energy    & 0.51\% & 1.18\% & -- \\
    Molecular and atomic nitrogen                    & 0.02\% & 0.13\% & -- \\
    Molecular and atomic nitrogen, with reactions    & 0.03\% & 0.18\% & 0.03\% \\
    \bottomrule
  \end{tabular}
\end{table}
